% Internal, AI-assisted research draft.  Author verification and the target
% journal's current AI policy must be resolved before submission.
% Bibliography keys are defined in ../references_additions.bib.
\subsection{関連する中心の問題と動機}
\label{subsec:related-centers}

物体の内部に特別な点を選ぶことには、変分的な観点からの相当な
歴史がある。Moszy\'nskaの動径中心は、選んだ点からの動径距離の関数の球面
平均を最大化することで定義され、
凸体に対しては、凹かつ狭義単調増加な平均化プロファイルが一意的な
最大点を与える\citep{moszynska2000selectors}。O'Haraは、Rieszポテンシャルと
その正則化から一般化された中心を展開している
\citep{ohara2012renormalization}。ポテンシャルの次数を$\beta$と書くと、
正則化していない式は次のとおりである：
\[
 V_{\Omega}^{(\beta)}(p)
   =\int_{\Omega}|x-p|^{\beta-n}\,dx,\qquad \beta>0.
\]
問題となる極値は次数に依存し、$\beta=n$では対数による置き換えが、
非正の次数では正則化が伴う。滑らかな凸体
$K\subset\mathbb R^3$と$0<\beta<3$に対して、このRieszポテンシャルの発散定理による
表示は次のとおりである：
\[
 V_{K}^{(\beta)}(p)
  =\frac{1}{\beta}\int_{\partial K}|x-p|^{\beta-3}
       (x-p)\cdot n(x)\,dA(x).
\]
幾何学的な因子$(x-p)\cdot n(x)\,dA(x)$は、まさに
以下で用いる基点に依存する錐体積の重みである。O'Haraのポテンシャルは
これに距離のべきを掛けて、極値を与える
中心を選ぶのに対し、本論文のエネルギーはこれに、動径方向と外向き法線の間の
角の二乗を掛け、
回転楕円体の軸比が変わるときのすべての停留成分を調べる。これらの構成は、
基点そのものを幾何学的汎関数の変数として
扱うことの自然な先例を与える。

これらの文献における一意性は、汎関数と物体の幾何の
両方に依存する。O'Haraは、与えられた各ポテンシャル次数において、非球面性が十分小さいとき、極値を与えるRieszポテンシャル中心の
一意性を証明している\citep{ohara2022nearly}。Sakataは、
狭義単調減少な動径核をもつポテンシャルの最大点を研究し、
一意性のための追加の十分条件を証明し、
最大点が二つ以上ある例を調べている\citep{sakata2016uniqueness}。とくに、平面の円環
$A=\{y\in\mathbb R^2:1\leq|y|\leq2\}$に対しては、その
$r^{-1/2}$ポテンシャル中心の集合は円であり、円板の和集合は
そのような中心を二つもちうる。これらの例は非凸な平面領域であるが、
本論文の円はこれに対し、凸な三次元の扁長
回転楕円体の内部における角度エネルギーのMorse--Bott\emph{鞍点}集合として
現れる。これらの例は、極値集合の構造に
注意深く目を向ける動機となるが、ここで考える角度
汎関数の停留集合を分類するものではない。また、選ばれた大域的
極値点の一意性は、それだけでは、別の
汎関数のすべての臨界点を記述するものでもない。

我々の問題は、$p$から境界点$x$へ向かう方向と
$x$における外向き法線との間の角の二乗を、次の
基点に依存する確率測度で平均したものに関わる：
\[
 d\mu_{K,p}(x)
   =\frac{(x-p)\cdot n(x)}{3\,\operatorname{Vol}(K)}\,dA(x)
   \qquad (p\in\operatorname{int}K).
\]
したがって、この定義には、視線
方向と重みに加えて法線場も関わる。我々は、このエネルギーを動径核の
体積ポテンシャルと同一視せず、これに対して前述の一意性定理を援用することもしない。
その代わりに、$p$に関する導関数が
消える集合の全体を調べる。回転楕円体族については、パラメータ$\lambda$は体積を固定して形状を
変えるものであり、Riesz核の次数$\beta$ではない。目的は、
中心の臨界点、軸上の対、および回転対称な
臨界円が、それらの退化と境界への接近も含めて、
一つの停留集合の図式にどのように収まるかを決定することである。

動径Gauss写像は、凸体の境界上で、動径方向に外向き
法線を対応させる。Gauss像
問題におけるその役割は、B\"or\"oczky, Lutwak, Yang, Zhang, Zhaoによって研究されている
\citep{boroczky2020gaussimage}。我々の汎関数にも、これと同じ
動径–法線の対応が現れるが、基点$p$は物体の内部を
動く。比較のために述べると、BertrandによるGauss曲率
指定問題の最適輸送による定式化では、（同論文の式~(6)において）外向き法線
$n$と動径方向$u$の間の角度コスト
$c(n,u)=-\log\cos d_{\mathbb S^2}(n,u)$が用いられる\citep{bertrand2016gauss}。ここでは
コストはその代わりに角の二乗$d_{\mathbb S^2}(n,u)^2$であり、これを
錐体積の確率測度に関して平均するが、その停留点を分類するためにGauss像の定理や
曲率指定の定理を援用することはしない。

熱流の中心は、関連するが別種の比較対象を与える。
MagnaniniとSakaguchiは、温度が
$u_t=\Delta u$を満たし、初期温度が$1$、境界温度が$0$である有界な伝導体を研究している
\citep{magnaniniSakaguchi2004hotspot}。ホットスポットとは$u(\,\cdot\,,t)$を最大にする点であり、
彼らの定常ホットスポットの問題は、同じ点がすべての正の時刻において
最大点であり続けるかどうかを問う。そこで述べられている凸な伝導体の問題では、
固定した各正の時刻において、空間的な臨界点はただ一つである。
ここで調べる条件は、固定した各形状における$d_p E_\lambda=0$であり、
極値点だけでなく鞍点も含む。その定義に
熱の時間発展は現れない。したがって、角度の
問題に鞍点円が現れることは、ホットスポットの一意性の結果から導かれることも、それと矛盾することもない。

非有界な形状の範囲にわたってそのような図式を確立するには、
解析的な簡約と数値的な検証とを結ぶものが必要である。この結びつきは、
計算機援用証明の確立された実践に属するものであり、
Tuckerによる常微分方程式の厳密な解析がその例であって、彼による精度保証付き
数値計算の解説およびG\'omez-Serranoの概説において説明されている
\citep{tucker2002rigorous,tucker2011validated,gomezSerrano2019survey}。
JohanssonのArbは、任意精度の
中心半径型演算の標準的な実装を提供し\citep{johansson2017arb}、彼の後の
論文は、ボール演算における厳密な数値積分を記述している
\citep{johansson2018integration}。演算による包囲
だけでは、特定の積分核の正しさや
領域被覆の完全性は確立されない。証明は、どの
不等式が排除、存在、一意性、および型を導くのか、また
有限被覆が解析的な端点領域とどのように接続するのかも、明示しなければならない。我々は、
精度保証付き数値計算の枠内での、公式、
包囲、および検査のこの問題固有の組織化を\emph{GeoEnclosure}と呼ぶ。


区間Newton判定とKrawczyk判定は、箱の内部の正則な根を分離できる\citep{intervalNewton}。分枝刈り込み型の解法は、局所的な判定と孤立根の大域的な探索とを組み合わせる\citep{branchPrune}。半径多項式法は、Newton型写像が縮小写像となる球を見つけることで、局所一意性を証明できる\citep{radiiContinuation}。停留円においては、全勾配のJacobianが軌道に沿って至るところ特異であり、軌道のどの点も孤立していない。回転によって軌道座標が取り除かれ、スカラーの分離が意味をもつようになる。ここでの問題固有の寄与は、デカルト座標での接方向の四階導関数から動径方向の交差を復元すること、および直接の微分が接触に出会うところで弦Hessianの差に切り替えることである。被覆と分離の原理そのものは標準的である。

\paragraph{このエネルギーに関する先行研究。}
著者による以前の球に関する原稿は、球に対する同じ
停留点・最小点の結果を与えているが\citep{furutaSphere}、本論文の
命題~\ref{prop:sphere}の証明は自己完結しており、
その原稿に依拠しない。Furutaによる以前の扁長側の局所
分岐の計算\citep{furutaProlate}は、中心における赤道方向の二次係数の零点を、Arbによる端点評価によって
$[4.72438340452113340672,4.72438340452113340673]$の中に
分離している。その軸比は本論文の$\lambda$である。横断方向の長さの単位を変えても、二次係数に正の正規化定数を掛けても、この零点は変わらない。したがって、局所的な閾値と円の発生は本論文に先行する。大域的な議論にどの端点計算が入るのかを特定するために、表~\ref{tab:critical}には、我々の窓全体の被覆から得られるより広い挟み込み区間を残してある。同じプロジェクトの三軸の記録は、局所的な開折、Hessianの退化曲線の交わり、および二つの軸上のMorse指数を扱っている。明示的に述べた有限個の入力C1--C6のもとで、ここでの寄与は、それらの局所的な分岐を、軸比の全域にわたる停留集合の分類（軸上の鞍点対と赤道上のMorse--Bott鞍点円の範囲、それらが境界に到達する軸比、およびそれ以外の停留成分の排除）へと拡張することである\citep{furutaProlate,furutaTriax}。
